\section{Conclusion}
This work establishes the static structural baseline of \textit{Informational Energetics}, a framework deriving the minimal architecture of persistence from control theory, information theory, and thermodynamics. We established that any system resisting entropic decay must contain all six irreducible pillars (Capacity, Map, Protocol, Governor, Toll, Margin), each demonstrated by its catastrophic failure mode when removed.

We tested whether IE and the six pillars are the minimal architecture for persistence by pushing IE to its most demanding domain: the quantum vacuum. Implementing the three existential requirements (Finiteness, Unitarity, and Causality) through the six-pillar architecture strictly constrained the viable vacuum substrate to the $E_8$ lattice projected onto a causal 4D manifold. This projection intrinsically yields the spacetime geometry including the Lorentzian metric signature (a geometric resolution of the Distler-Garibaldi no-go theorem (\cref{app:geometric_origin_of_chiral_fermions})), temporal cycle ($\Delta = 43$), channel depth ($\nu = 16$), embedding dimension ($\sigma = 5$), and boundary topology ($\chi = 2$) demanded by the persistence architecture as mathematical consequences, with no free parameters. The unique isolation of a single substrate constitutes the first structural test.

The second, decisive test evaluates whether this substrate's geometric impedance $Z_{\text{geo}}$ (a structural index derived from the six-pillar entropic balance) correctly corresponds to the static, infrared fine-structure constant ($\alpha^{-1}(0) = \num{\AlphaInvVal}\ldots$). This value structurally aligns with the kinematic measurement of Morel (2020) at $\AlphaInvMorelSigma\sigma$ and CODATA (2022) at $\AlphaInvCodataSigma\sigma$, while predicting a strict \textit{Quantization Gap} where continuous, high-loop QED extractions must systematically diverge due to the substrate's finite channel capacity. The same six-pillar impedance determines the von Klitzing constant $R_K$ and the elementary charge $e = q_P/\sqrt{Z_{geo}}$, which emerges strictly as a flow constraint enforcing the Schwinger vacuum breakdown limit.

This result supplies what complex systems research has lacked: a formal boundary condition for persistence. Current frameworks model how systems self-organize through descriptive phenomena (adaptation, modularity, resilience) without deriving the minimal architecture that makes any of them possible. IE establishes that persistence is a single universal constraint. The six pillars constitute an eliminative filter: any system lacking one inevitably dissolves; any system possessing all six contains the necessary mechanics to resist entropic decay, regardless of substrate. The vacuum derivation is the most demanding test case. If the same architectural requirements constrain biological cells, computational networks, and quantum substrates alike, then persistence is substrate-universal, and complex systems models now possess a falsifiable, first-principles foundation.

This paper establishes the static architectural baseline. Natural avenues for the validation and extension of IE include defining the full dynamics of the Entropic Action and recursive partitioning into subsystems when internal complexity exceeds the capacity of a single control loop. Because any persistent subsystem must re-implement the same six-pillar architecture against its local entropic boundary, the framework naturally generalizes to nested, hierarchical scales. Having established the static baseline of the $E_8$ substrate, it serves as a derivational platform: these extensions provide an explicit roadmap to verify the emergence of General Relativity, the Standard Model Lagrangian, and its remaining dimensionless constants, which can then be compared against experimental values. Ultimately, the definitive proof that $\alpha^{-1}$ is an architectural necessity lies in whether this exact, unadjusted substrate generates the full Standard Model action and remaining constants when extended beyond quiescent equilibrium.

Informational Energetics proposes that persistence is the primary organizing principle of reality at every scale. If the fundamental vacuum is uniquely determined by the strict constraints of persistence, then by the Selection Principle, any persistent macroscopic system is solving the exact same architectural problem, differing only in substrate material and environmental openness. Ultimately, physical reality is the cleanest, most rigorous environment in which to first test the architecture of persistence.

\begin{acknowledgments}
The author is an independent researcher and received no external funding for this work. I would like to thank Brian Sheppard for rigorous and constructive feedback. I would like to thank my friends and family for their patience and support throughout decades of discussions, as I tried to understand every field I became interested in and the recurring, overarching patterns I saw.
\end{acknowledgments}